%% APPENDIX E -- radio-frequency interference screening.
%% Owner: appendix-triage.  Carries \label{app:rfi}.
\section{Radio-frequency interference screening}
\label{app:rfi}

Interference is the first alternative explanation for an unattributed
crossing, so the screening this survey has and the coverage it lacks are both
set out here. Millimetre interferometry is an unusually poor
environment for terrestrial interference, though not an empty one: ITU Radio
Regulations Article~5 \citep{ITU2024} carries active allocations throughout
90--275\,GHz, and WRC-19 identified parts of 275--450\,GHz for land-mobile and
fixed use. The specific hazard is the \RadarLo--\RadarHi\,GHz Earth
exploration-satellite allocation used by spaceborne cloud radar
\citep{CloudSat2002,EarthCARE2023}, a documented contaminant of ALMA Band~3
\citep{ALMAMemo504}: of the \RadarNTot{} census windows \RadarNWin{} overlaps
it, carrying \RadarNCross{} threshold crossings and \RadarNStage{}
rank-passing events, so no result here rests on it. Two things favour this
band: terrestrial transmitters are sparse compared with centimetre
wavelengths, and an interferometer is not a total-power detector, so a signal
not in the far field at the phase centre acquires a rapidly varying
geometric phase across the array and is suppressed in the correlation rather
than imaged.

\emph{The coverage of the screening is uneven.} The five tests of
Table~\ref{tab:rfitests} are computed for every searched window of the
census, as is the \NCtrl-position rank screen; recurrence
reaches only the crossings with a covering repeat block, and the visibility
fit is a consistency check rather than a filter. The correlator flagging is
the observatory's own, replayed from each block's delivered calibration
script, and was not retained window by window: \emph{no per-channel
interference flag of our own exists, and none can be recovered.} The
populations of the five tests in Table~\ref{tab:rfitests} differ from one
another for that reason: each test needs a different quantity to have been
recorded.

\begin{table}
\centering
\footnotesize
%% ★ v4.26 (r25-captions): 4 printed lines -> the author's 3.  The reason the
%% five populations differ MOVED to the paragraph above, where the record
%% that is or is not retained is the subject.
\caption{The five interference tests, each of which could have failed. $p$ is
against the same peaks redrawn on their own channel grids.}
\label{tab:rfitests}
\begin{tabular}{@{}p{0.20\columnwidth}p{0.22\columnwidth}p{0.33\columnwidth}r@{}}
\hline
Test & Population & Result & $p$ \\
\hline
Clustering in sky frequency &
\AppMNCrossF{} crossings with a measured peak &
\AppMTubeObs{} within $\pm\DispoHalfKms$\,km\,s$^{-1}$ of a searched
transition against \AppMTubeNull{} expected: they cluster on the survey's own
lines &
\AppMTubeP \\
Crowding on one another &
\QcNQP{} unattributed, quality-passing &
no crowding at any of six separations from \QcTolLo{} to \QcTolHi\,MHz, where
the \QcNAttr{} attributed crossings crowd at every one &
$\geq\QcPQPMin$ \\
Atmospheric transitions &
\TlNCross{} crossings against \TlNLine{} telluric lines &
\ScNInRad{} inside the \ScRadMHz\,MHz screening radius in the topocentric
frame, flagged as atmospheric and removed from both sides of the chance
comparison while remaining among the unattributed; \TlTubeObs{} inside
$\pm\TlTubeHalf$\,km\,s$^{-1}$ against \TlTubeNull{} expected &
-- \\
Fixed intermediate frequency &
\IfPopNWin{} windows with a released peak frequency &
peak positions within their own windows uniform, median \RfiFracMed &
-- \\
Cross-target occupancy &
\RfiOccNWin{} windows toward \RfiOccNStars{} stars &
\RfiOccObs{} cross-star coincidences within \RfiOccTol\,MHz against
\RfiOccNull$\,\pm\,$\RfiOccNullSd{} expected; \emph{below} the null over the
\RfiOccSigN{} above-trigger windows, \RfiOccSigObs{} against
\RfiOccSigNull$\,\pm\,$\RfiOccSigNullSd &
\RfiOccP \\
Image sideband &
\ImNScreen{} unattributed crossings &
\ImNNoParent{} have no catalogued transition in the image interval and
\ImNParentAbsent{} have one the same block observed with nothing bright
enough to leak; for \ImNLeft{} the parent frequency was not observed &
-- \\
\hline
\end{tabular}
\end{table}

The first four tests establish that the crossing population carries none of
the signatures interference leaves. The clustering it does show is the
attribution and not a residue left over after it: the ledger counts the
densest band, \AppMBandLo--\AppMBandHi\,GHz, as \EvTubeLedgerN{} crossings
with \EvTubeLedgerAttr{} attributed, and the one concentration a reader will
notice, \TlBandAll{} crossings between \TlBandLo{} and \TlBandHi\,GHz, is the
carbon monoxide tuning seen from the other side -- \TlBandObs{} of the
\TlBandNSub{} unattributed crossings whose window the catalogue tabulates
against \TlBandNull$\,\pm\,$\TlBandNullSd{} expected from where the spectral
windows were tuned. It is not atmospheric either, because that interval
holds \TlBandNLine{} telluric transitions. The mask has to be applied before the
occupancy screen, because \RfiOccMaskN{} of the \RfiOccSigN{} above-trigger
windows sit on a masked transition (\RfiOccSpeciesList), and a carbon
monoxide line falls at nearly the same sky frequency toward every nearby
star, which is exactly the coincidence that screen is looking for; applied,
it reads \RfiOccMskObs{} against
\RfiOccMskNull$\,\pm\,$\RfiOccMskNullSd. The atmospheric screen is not sharp
where it fires: the one crossing it removes sits in a window tuned on the
line, \ScNearWinPct{} per cent of whose channels lie within the same radius
of the same transition, so the coincidence establishes not that this channel
is atmospheric but that this window cannot separate the two, and removing it
from both sides of the comparison of \S\ref{sec:trigunif} is the
conservative response. It remains among the \MkUnattr{} unattributed
crossings and is tested for recurrence with them. The same radius is taken off the control side,
where it enters the expectation of Appendix~\ref{app:falsealarm} as one
surviving fraction formed over all windows rather than crossing by crossing;
that pooled fraction is the larger of the two, so the screen takes
proportionally more from the expectation than from the observed count and the
residual asymmetry runs against the null and not for it.

\emph{The fifth test bounds the one mechanism that could put a crossing where
no line is catalogued.} A sideband-separating receiver passes the
image frequency $2\nu_{\rm LO1}-\nu$ at finite rejection, and the first
local oscillator is not among the quantities the archive delivers; each
block's own spectral windows bound it, every window having to
fall inside one sideband at an intermediate frequency the receiver passes,
which pins $\nu_{\rm LO1}$ to a median \ImLoWidthBlkGHz\,GHz. Reflecting each
crossing about that interval places its image \ImGapLoGHz--\ImGapHiGHz\,GHz
away, never inside the window the crossing was found in, and ALMA specifies
an image rejection of at least \ImIrrDb\,dB, so a parent would have to be
\ImIrrFac{} times the crossing's own flux density. Taken \ImWideGHz\,GHz
wider at each end the screen leaves \ImNNoParentWide{} of the \ImNScreen{}
excluded on frequency alone, so it is a bound and not a clearance, and
\ImNDsb{} unattributed crossings lie in Band~9 or Band~10, where the
receivers are double-sideband and no such bound could be drawn. Where a
carbon monoxide isotopologue does fall in the image band of a Class~A window,
\ScImNHitCross{} crossings lie at the image position against \ScImExp{}
expected from the share of the window that band occupies, \ScImNHitUn{} of
them unattributed -- and neither parent attributes its crossing: each was
observed at its own frequency in the same block, where no
feature reaches the trigger, \ScImUnParMaxT{} at most, and the largest
amplitude is \ScImUnRatioA{} and \ScImUnRatioB{} times the crossing's own,
where a leak at the quoted rejection would need \ImIrrFac{} times as much.

\emph{One execution block fails these tests, and it fails them as a block
and not as a carrier.} Its \DqBlockNCross{} crossings toward
\LgVerRejectStar{} are the \DqBlockNCross{} highest statistics in the survey
(\S\ref{sec:bdblock}, which identifies the failure). Whether interference
accounts for them is what belongs here, and three readings say not. The
visibility fit returns a large imaginary part,
placing the emission away from the phase centre; the block's
\BdNSibling{} sibling executions of the same field, array and spectral setup
produce no crossing at all; and of the further windows covering each of its
frequencies toward other stars, \DqOtherSame{} carries a crossing in the
same channel.

\emph{None of this substitutes for the ON/OFF cadence a dedicated search
uses.} The claim is only that no
identified interference mechanism accounts for the unattributed events:
allocation tables cannot exclude harmonics, intermodulation products,
unconventional emitters or near-field sources, and archival data without a
simultaneous off-source strategy cannot do better. These are statements about
the population, not a clearance of any one crossing.

%% "The null result rests on internal validation alone" lives once, in the
%% discussion's lessons subsection, and not here.  fig:appm is not included
%% -- it compared our own two statistics and the paper states only the one it
%% uses -- so figures/appm_power.pdf and the \AppMPowerOld* and
%% \RfiOcc{Prior,Dir,Stale}* families will be retired.
%%
%% ★ r16-mask: the cross-target screen of the fourth test is now scoped to
%% the ADOPTED attribution (round 660, maskrfi_v660.py), which makes it two
%% events and not one.  round 49's \RfiEventLead, \RfiFreqThem,
%% \RfiSameChanClause, \RfiNotTestableClause, \RfiOtherWin,
%% \RfiOtherStarsRange and \RfiOtherBlocksRange are therefore unreferenced
%% and will be retired; rfi_v399.py still reproduces them and its grammar
%% selftest still runs, so nothing breaks if the scope is ever restored.
